\documentclass[11pt,a4paper]{article}

\usepackage[utf8]{inputenc}

\usepackage[T1]{fontenc}

\usepackage[polish]{babel}

\usepackage{lmodern}

\usepackage{geometry}

\geometry{margin=27mm}

\usepackage[hidelinks]{hyperref}

\title{Rozbieżność aktualizacji faktów i norm moralnych}

\author{Mateusz Pomianek}

\date{Szkic naukowy, wersja 1.0.0 | 28 listopada 2026}

\begin{document}

\maketitle

\begin{abstract}

Proponowana miara główna to trójczynnikowa interakcja warunku korekty, czasu i rodzaju oceny.

\end{abstract}

\section{Wprowadzenie}

Korekty faktów mogą poprawiać trafność przekonań (Wood \& Porter, 2019). Oceny moralne uwzględniają również intencje i odpowiedzialność (Cushman, 2008), dlatego podobna dynamika obu ocen wymaga odrębnego testu.

\section{Cel i pytania badawcze}

Hipotezy robocze wymagają zatwierdzenia przed rozpoczęciem właściwej weryfikacji.

\paragraph{H1.} Korekta dostarczy większej zmiany oceny zgodnej z informacją eksperymentalną w przypadku twierdzeń o faktach niż w przypadku oceny normatywnej powiązanej z tym samym scenariuszem, względem warunku neutralnego.

\section{Materiał i metody}

Pierwszorzędowy wskaźnik: trójczynnikowa interakcja warunku korekty, czasu i rodzaju oceny.

Fakt i normę należy mierzyć przed komunikatem i po nim na skalach o tej samej rozpiętości, a bodźce dobierać tak, aby przed badaniem ustalić oczekiwany kierunek obu zmian. Główny kontrast to interakcja czasu, rodzaju oceny i warunku komunikatu. Wiarygodność źródła można zróżnicować niezależnie.

Przed badaniem z udziałem ludzi należy zatwierdzić wymagania etyczne, zasady zgody i ochronę danych.

\section{Plan analizy}

Zamrożenie estymandu, reguł oceny jakości i obsługi braków nastąpi przed dostępem do danych wynikowych. Raport obejmie oszacowanie kontrastu głównego z przedziałem niepewności i pełne mianowniki. Jeśli przewidziano hipotezy drugorzędowe, przed analizą należy ustalić sposób kontroli wielokrotności.

\section{Wyniki}

W niniejszym pliku nie umieszczono wyników empirycznych, statystyk ani nieprzeprowadzonych analiz. Sekcja wymaga uzupełnienia po uzyskaniu i weryfikacji materiału dowodowego.

\section{Dyskusja i ograniczenia}

Ograniczenia obejmują możliwe błędy pomiaru, dobór próby, alternatywne mechanizmy i ograniczoną możliwość uogólniania. Korelacja nie dowodzi przyczynowości; poprawność techniczna nie jest dowodem znaczenia klinicznego.

\section{Wnioski}

Hipotezy i sposób pomiaru stanowią podstawę oceny wykonalności oraz przygotowania właściwego protokołu.

\section*{Bibliografia}

{\setlength{\parindent}{0pt}
\setlength{\parskip}{5pt}
\hangindent=1.27cm
Cushman, F. (2008). Crime and punishment: Distinguishing the roles of causal and intentional analyses in moral judgment. \textit{Cognition, 108}(2), 353--380. \url{https://doi.org/10.1016/j.cognition.2008.03.006}\par
\hangindent=1.27cm
Wood, T., \& Porter, E. (2019). The elusive backfire effect: Mass attitudes’ steadfast factual adherence. \textit{Political Behavior, 41}(1), 135--163. \url{https://doi.org/10.1007/s11109-018-9443-y}\par
}

\end{document}
